\documentclass[conference,a4paper]{IEEEtran}
\IEEEoverridecommandlockouts

\usepackage{cite}
\usepackage{amsmath,amssymb,amsfonts}
\usepackage{graphicx}
\usepackage{textcomp}
\usepackage{xcolor}

\begin{document}

% ============================================================
%  TITEL
% ============================================================
\title{Digital Twins for Predictive Maintenance in Automotive Aftersales}

% ============================================================
%  AUTOREN  -- alle 3 Gruppenmitglieder eintragen
% ============================================================
\author{
  \IEEEauthorblockN{Oliver Feng}
  \IEEEauthorblockA{
    \textit{Wirtschaftsinformatik} \\
    \textit{DHBW Stuttgart} \\
    Stuttgart, Germany \\
    wi23030@lehre.dhbw-stuttgart.de
  }
  \and
  \IEEEauthorblockN{Maksim Bogachenkov}
  \IEEEauthorblockA{
    \textit{Wirtschaftsinformatik} \\
    \textit{DHBW Stuttgart} \\
    Stuttgart, Germany \\
    wi23024@lehre.dhbw-stuttgart.de
  }
  \and
  \IEEEauthorblockN{Joel Yerai Martinez Campo}
  \IEEEauthorblockA{
    \textit{Wirtschaftsinformatik} \\
    \textit{DHBW Stuttgart} \\
    Stuttgart, Germany \\
    wi23194@lehre.dhbw-stuttgart.de
  }
}

\maketitle

% ============================================================
%  ABSTRACT  (keine Formeln, Sonderzeichen oder Symbole!)
%  Ziel: verstaendlicher Ueberblick ueber Thema, Vorgehen,
%        wichtigste Ergebnisse und Fazit (~150-250 Woerter)
% ============================================================
\begin{abstract}
TODO: Abstract here.
\end{abstract}

\begin{IEEEkeywords}
Digital Twin, Predictive Maintenance, Industry~4.0, Machine Learning, IoT, Condition Monitoring
\end{IEEEkeywords}


% ============================================================
\section{Introduction}
% ============================================================

\begin{itemize}
  \item Hook
  \item Smart to Wise Bezug
  \item Problemstellung
  \item Zeilsetzung
  \item RQ, Forschunslücke
  \item Beitrag der Arbeit
  \item Aufbau der Arbeit
\end{itemize}


% ============================================================
\section{Related Works \& Theoretical Background}
% ============================================================

\subsection{Digital Twins}

\subsection{Vehicle Maintenance Strategies}

Vehicle maintenance strategies form a maturity hierarchy: reactive maintenance repairs only after failure \cite{abougrad2026xai}; preventive maintenance schedules service at fixed mileage or time intervals, leading to either premature replacement or undetected wear \cite{abougrad2026xai,aichi2024multilayered}; condition-based maintenance (CBM) couples the service decision to sensor thresholds \cite{aichi2024multilayered}; and predictive maintenance (PdM) extends CBM by forecasting the future evolution of degradation, estimating the remaining time or distance until end-of-life is reached \cite{kumar2025suspension,abougrad2026xai}.

PdM relies on three data layers: on-board sensors delivering high-frequency signals such as vibration, current, voltage, and cell temperature \cite{kumar2025suspension,karthick2025simulating,aichi2024multilayered}; the On-Board Diagnostic (OBD) interface providing standardized diagnostic codes \cite{zhang2009connected,coppola2016connected}; and over-the-air (OTA) telematics enabling continuous low-latency exchange with the backend \cite{lopes2025vehicledt,karthick2025simulating,coppola2016connected}. Together, these layers transform the vehicle from a periodically inspected asset into a continuously observed system, a shift first documented by Zhang et al.~\cite{zhang2009connected} in their work on Connected Vehicle Diagnostics and Prognostics at General Motors.

PdM at this scale is not a stand-alone technique but a service layer built on a digital twin (DT) of the vehicle. As Lopes et al.~\cite{lopes2025vehicledt} and Hadraoui et al.~\cite{aichi2024multilayered} describe, the DT supplies the synchronized virtual representation against which sensor streams are interpreted, anomalies detected, and remaining-useful-life estimates computed. Without a DT, telemetry must be processed in isolation; without PdM, the DT remains a passive monitor. In current automotive DT architectures, predictive analytics is therefore typically positioned as a dedicated microservice within the DT's analytics core \cite{lopes2025vehicledt,karthick2025simulating,abougrad2026xai}.

\subsection{Machine Learning in PdM}

Three problem classes structure ML-based PdM. Classification assigns sensor patterns to discrete fault categories using Support Vector Machines, Random Forests, or Convolutional Neural Networks on time-frequency representations of vibration signals \cite{kumar2025suspension,abougrad2026xai}. Regression targets Remaining Useful Life (RUL) estimation, dominated by Long Short-Term Memory (LSTM) networks for their ability to model long-range temporal dependencies in degradation signals \cite{karthick2025simulating,abougrad2026xai}. Anomaly detection operates without explicit labels, learning nominal behavior and flagging deviations as failure precursors \cite{karthick2025simulating,abougrad2026xai}. AbouGrad et al.~\cite{abougrad2026xai} report classification accuracies of approximately 89\% for vibration-based fault detection and 92\% for LSTM-based RUL estimation in automotive contexts, though they emphasize that these figures depend strongly on dataset quality and operating conditions.

A recurring constraint is data: AbouGrad et al.~\cite{abougrad2026xai} and Hadraoui et al.~\cite{aichi2024multilayered} note that preprocessing consumes 70--90\% of project time, genuine failure data is structurally scarce, and public benchmarks representative of in-service behavior are largely absent. Model-based and hybrid approaches that embed physical degradation laws into the learning pipeline are therefore preferred over purely data-driven ones, as they extrapolate more reliably under data scarcity \cite{kumar2025suspension,aichi2024multilayered,wang2023review}. Beyond accuracy, deployment in safety-relevant contexts raises the question of explainability: deep architectures achieve high accuracy but operate as opaque models whose decisions cannot readily be justified to engineers, customers, or regulators \cite{abougrad2026xai}. Explainable AI (XAI) techniques such as feature attribution and attention-based saliency are increasingly integrated into PdM pipelines to expose the reasoning behind a prediction, marking the conceptual bridge from smart to wise information systems \cite{abougrad2026xai}.

\subsection{Automotive Aftersales \& Customer Lifecycle}


\texttt{Forschungslücke}

While these contributions establish the technical foundations of DT-based PdM, existing work concentrates either on algorithmic aspects in isolation \cite{kumar2025suspension,abougrad2026xai} or on industrial and aerospace applications \cite{wang2023review}. The integration of DT-based PdM into automotive aftersales business models --- including its implications for customer interaction, upselling potential, and vehicle lifecycle value --- remains underexplored. The present work addresses this gap.
% ============================================================
\section{Methodology}
% ============================================================

\begin{itemize}
  \item LR
  \item Ilustrative Case Study
  \item (DSR)
\end{itemize}


% ============================================================
\section{Reference Architecture}
% ============================================================

\begin{itemize}
  \item Sensorik
  \item ML-Algorithmen
  \item Kennzahlen
  \item Schnittstellen
  \item Datenmanagement
\end{itemize}


% ============================================================
\section{Ilustrative Case Study}
% ============================================================

\subsection{Lifecycle Comparison}

\subsection{Assumption Framework}

\subsection{Ilustrative Calculation \& Results}

\subsection{Limitations}


% ============================================================
\section{Discussion and Recommendations}
% ============================================================

\begin{itemize}
  \item Interpretation der Ergebnisse
  \item Implikationen fuer Forschung und Praxis
\end{itemize}


% ============================================================
\section{Conclusion \& Future Work}
% ============================================================

\begin{itemize}
  \item Zusammenfassung der wichtigsten Erkenntnisse
  \item Beitrag der Arbeit
  \item Ausblick auf weitere Forschung \cite{grieves2014}
\end{itemize}


% ============================================================
\section*{Acknowledgment}
% ============================================================
The authors would like to thank Eduard Dobenko (Mercedes-Benz)
for supervising this work.


% Quellen in references.bib eintragen, hier nichts aendern
\bibliographystyle{IEEEtran}
\bibliography{references}


% ============================================================
%  KI-ERKLAERUNG
%  Laut DHBW-Vorgabe auf einer EIGENEN SEITE am Ende,
%  damit sie vor der Einreichung leicht abgetrennt werden kann.
%  Diese Seite erscheint NICHT im Sammelband.
% ============================================================
\newpage
\section*{Erklärung zur Nutzung von KI-Tools}

TODO: KI-Erklaerung hier einfuegen (gemaess DHBW-Vorgabe).

\end{document}
